\subsubsection{What the device model relies on}

For the final two-foot model (person scaling, forest vote, 92.3\,\%), shuffling all features of one foot within each left-out
participant costs 33.4 points (trailing foot) and 24.5 points (leading foot): the model uses both feet
(Table~\ref{tab:importance}). Within a foot, the foot-trajectory summaries matter most (about 6 points each), followed by the
clearance-shape features (about 3) and the foot tilt and angular velocity (about 2); the foot acceleration hardly matters. No single
feature costs more than 2.5 points (Table~\ref{tab:importancetop}): the information is spread over redundant features, the most
useful being how high each foot rises (peak foot height, height at 25\,\% of the trailing swing) and the trailing stance time.

\begin{table}[htbp]
\centering
\begin{minipage}{0.52\textwidth}
\centering
\caption{Permutation importance of feature groups (accuracy drop, mean $\pm$ SD over 10 shuffles).}
\label{tab:importance}
\begin{adjustbox}{max width=\textwidth}
{\footnotesize
\begin{tabular}{lrr}
\toprule
\textbf{shuffled feature group} & \textbf{features} & \textbf{accuracy drop (points)} \\
\midrule
trail foot, all features & 41 & 33.4 $\pm$ 3.4 \\
lead foot, all features & 46 & 24.5 $\pm$ 1.7 \\
trail: foot trajectory summaries & 11 & 6.0 $\pm$ 0.9 \\
lead: foot trajectory summaries & 11 & 5.7 $\pm$ 1.0 \\
trail: clearance shape & 13 & 2.9 $\pm$ 0.8 \\
lead: clearance shape & 13 & 2.9 $\pm$ 1.0 \\
trail: foot tilt and angular velocity & 14 & 2.2 $\pm$ 1.0 \\
lead: foot tilt and angular velocity & 14 & 2.2 $\pm$ 1.0 \\
trail: timing and approach strides & 2 & 1.9 $\pm$ 0.7 \\
lead: timing and approach strides & 7 & 0.7 $\pm$ 0.5 \\
trail: foot acceleration & 3 & 0.7 $\pm$ 0.3 \\
lead: foot acceleration & 3 & 0.2 $\pm$ 0.3 \\
\bottomrule
\end{tabular}
}
\end{adjustbox}

\end{minipage}\hfill
\begin{minipage}{0.44\textwidth}
\centering
\caption{The 12 most important single features.}
\label{tab:importancetop}
\begin{adjustbox}{max width=\textwidth}
{\footnotesize
\begin{tabular}{lr}
\toprule
\textbf{feature} & \textbf{accuracy drop (points)} \\
\midrule
\texttt{trail: trail\_\allowbreak{}zhc\_\allowbreak{}peak\_\allowbreak{}foot\_\allowbreak{}h} & 2.5 $\pm$ 0.7 \\
\texttt{trail: trail\_\allowbreak{}zhc\_\allowbreak{}h\_\allowbreak{}25} & 2.5 $\pm$ 0.8 \\
\texttt{trail: shape\_\allowbreak{}h\_\allowbreak{}peak} & 2.5 $\pm$ 1.0 \\
\texttt{lead: shape\_\allowbreak{}h\_\allowbreak{}peak} & 2.1 $\pm$ 0.9 \\
\texttt{lead: lead\_\allowbreak{}zhc\_\allowbreak{}peak\_\allowbreak{}foot\_\allowbreak{}h} & 2.0 $\pm$ 1.0 \\
\texttt{trail: trail\_\allowbreak{}stance\_\allowbreak{}time} & 2.0 $\pm$ 0.7 \\
\texttt{trail: trail\_\allowbreak{}zhc\_\allowbreak{}h\_\allowbreak{}toeoff} & 2.0 $\pm$ 0.8 \\
\texttt{trail: trail\_\allowbreak{}zhc\_\allowbreak{}h\_\allowbreak{}50} & 1.1 $\pm$ 0.6 \\
\texttt{lead: lead\_\allowbreak{}zhc\_\allowbreak{}h\_\allowbreak{}50} & 0.9 $\pm$ 0.6 \\
\texttt{trail: shape\_\allowbreak{}rise\_\allowbreak{}fall} & 0.9 $\pm$ 0.4 \\
\texttt{lead: cadence} & 0.7 $\pm$ 0.6 \\
\texttt{trail: trail\_\allowbreak{}foot\_\allowbreak{}av\_\allowbreak{}min} & 0.6 $\pm$ 0.3 \\
\bottomrule
\end{tabular}
}
\end{adjustbox}

\end{minipage}
\end{table}

\subsubsection{Statistics for the main comparisons}

\paragraph{Paired differences.} On the same crossings (Table~\ref{tab:paired}), person scaling adds 34.8 points to the two-foot
model (95\,\% CI +22.3 to +46.2), and the 13 clearance-shape features add 10.7 points to a generic model (+4.9 to +16.9). Six IMUs
are not better than two foot IMUs (+1.1, CI $-0.5$ to +3.0), and the trailing foot adds 3.5 points to the leading foot alone (CI
$-1.0$ to +7.2, not significant with ten participants).

\begin{table}[htbp]
\centering
\caption{Paired comparisons on the same crossings (participant bootstrap of the difference).}
\label{tab:paired}
\begin{adjustbox}{max width=\textwidth}
{\footnotesize
\begin{tabular}{p{7.8cm}rrr}
\toprule
\textbf{comparison} & \textbf{A} & \textbf{B} & \textbf{A - B (95\% CI)} \\
\midrule
person scaling vs generic (two foot IMUs, forest) & 92.3\% & 57.5\% & +34.8 (+22.3 to +46.2) \\
two foot IMUs vs all six IMUs (person scaling, forest) & 92.3\% & 91.2\% & +1.1 ($-$0.5 to +3.0) \\
both feet vs leading foot only (person scaling, forest) & 92.3\% & 88.8\% & +3.5 ($-$1.0 to +7.2) \\
13 clearance-shape vs 57 summary features (generic, logistic, leading limb) & 72.8\% & 62.1\% & +10.7 (+4.9 to +16.9) \\
\bottomrule
\end{tabular}
}
\end{adjustbox}

\end{table}

\paragraph{Errors.} The final model recognises 95\,\% of H1, 90\,\% of H2 and 91\,\% of H3 crossings; its errors are almost all
between neighbouring heights (one H3 crossing was called H1; Table~\ref{tab:confusion}).

\paragraph{Can the probabilities be trusted?} The person-scaled forest is \emph{under}-confident (Table~\ref{tab:reliability},
Figure~\ref{fig:validate}): crossings for which it states 55\,\% are right 80\,\% of the time, and from 0.7 upward it is right
96--100\,\% of the time. Its Brier score is 0.18, against 0.51 for the generic model (Table~\ref{tab:brier}). A confidence
threshold is therefore safe for this model: the stated probability understates its reliability.

\begin{table}[htbp]
\centering
\begin{minipage}{0.4\textwidth}
\centering
\caption{Confusion matrix of the final model.}
\label{tab:confusion}
\begin{adjustbox}{max width=\textwidth}
{\footnotesize
\begin{tabular}{lrrrr}
\toprule
\textbf{true \textbackslash{} predicted} & \textbf{H1} & \textbf{H2} & \textbf{H3} & \textbf{correct} \\
\midrule
H1 & 84 & 4 & 0 & 95\% \\
H2 & 0 & 81 & 9 & 90\% \\
H3 & 1 & 6 & 75 & 91\% \\
\bottomrule
\end{tabular}
}
\end{adjustbox}

\end{minipage}\hfill
\begin{minipage}{0.56\textwidth}
\centering
\caption{Brier score and expected calibration error (ECE).}
\label{tab:brier}
\begin{adjustbox}{max width=\textwidth}
{\footnotesize
\begin{tabular}{lrr}
\toprule
\textbf{model} & \textbf{Brier} & \textbf{ECE} \\
\midrule
final (person scaling) & 0.182 & 16.6 points \\
generic (level-walk scaling) & 0.506 & 9.4 points \\
\bottomrule
\end{tabular}
}
\end{adjustbox}

\end{minipage}
\end{table}

\begin{table}[htbp]
\centering
\caption{Reliability: crossings grouped by the stated top-class probability; number, mean stated probability and observed accuracy.}
\label{tab:reliability}
\begin{adjustbox}{max width=\textwidth}
{\footnotesize
\begin{tabular}{lrr}
\toprule
\textbf{stated probability} & \textbf{final: crossings, mean stated, observed} & \textbf{generic: crossings, mean stated, observed} \\
\midrule
0.33--0.50 & 15, 48\%, 67\% & 85, 44\%, 44\% \\
0.50--0.60 & 40, 55\%, 80\% & 53, 54\%, 49\% \\
0.60--0.70 & 30, 65\%, 90\% & 38, 65\%, 32\% \\
0.70--0.80 & 55, 75\%, 96\% & 39, 75\%, 82\% \\
0.80--0.90 & 68, 85\%, 97\% & 39, 84\%, 90\% \\
0.90--1.00 & 52, 94\%, 100\% & 6, 93\%, 100\% \\
\bottomrule
\end{tabular}
}
\end{adjustbox}

\end{table}

\begin{figure}[htbp]
\centering
\includegraphics[width=\textwidth]{fig_validate.png}
\caption{Validation. V1: height error when a height level is missing from training (grey = trained on all heights, green = this
height left out; light = ridge, dark = forest). V3: what explains the crossing-to-crossing clearance (blue = leading, orange =
trailing limb). V4: accuracy lost when a feature group of the final model is shuffled (permutation importance). V5: reliability of
the stated probabilities (grey = generic, purple = person scaling, dashed = perfect calibration).}
\label{fig:validate}
\end{figure}
